% Generated by task tex-opening from dossiers/gt-opening-compile.md (section H). Every probe
% source and output is the dossier's.
% The citation check dossiers/verify-gt-opening.md: see the note at the top of index-opening.tex.
\section{The opening and the compilation: probes}\label{sec:gen-opening-probes}

\begin{sloppypar}
All files are under \file{.claude/reports/blueprint-review/probes/gt-opening-compile/}. The two Lean
probes were run from the repository root with
\code{flock .claude/reports/blueprint-review/logs/lean.lock lake env lean <file>}, first on
2026-09-29 against the old pins (leanerVM \code{b435631}, ArkLib \code{dca90385}, CompPoly
\code{3468b38c}, VCVio \code{f9dc47d9}; both \code{exit=0}), then again on 2026-09-30 after the
rebuild, against the new pins (ArkLib \code{7653a901}, CompPoly \code{572f9973}, VCVio
\code{a4232d08}, Lean 4.34.1; both \code{exit=0}, log \file{rerun-newpin.log}). They import only
modules the build of \code{main} compiles; ArkLib's \file{Commitments/Functional/Basic.lean} and
\file{FiatShamir/Basic.lean} are \textbf{not} compiled by that build (no leanerVM module imports
them: 47 ArkLib \code{.olean}s under \file{.lake/packages/Arklib/.lake/build} on 2026-09-29, none
under \file{Commitments/}, \file{FiatShamir/}, \file{Data/CodingTheory/}), so the two definitions the
probes need from them are copied verbatim, with their source lines named, and the copies are what
elaborated. The two Python scripts are scratch computations; they run no Lean.
\end{sloppypar}

\subsection{\texttt{InnerProductOracle.lean}: the inner-product interface typechecks, and the opening phase is one query}\label{sec:gen-opening-probe-inner}

\paragraph{What it establishes.} That the inner-product oracle interface recommended in
\cref{sec:gen-opening-alternatives} is expressible at the pins: (1) ArkLib's \lean{OracleInterface}
accepts \lean{Weight n} (leanerVM's structure, in \code{Type}) as a query type; (2) the evaluation
query is the special case of the equality weight, by the built lemma \lean{eqWeight_pair}; (3)
ArkLib's \lean{Commitment.Scheme} over that interface has as opening statement (commitment, weight,
claimed value): an inner-product commitment scheme; (4) the opening phase ``one challenge $λ$, then
one weighted query'' is an ArkLib \lean{OracleVerifier} over a one-round schedule. A wrapper
\lean{Stack} is used so that the probe's instance does not overlap the built \lean{evalOracle}.

\paragraph{Source.}
\begin{leancode}
/-
Probe (task gt-opening-compile): is alternative (ii) expressible at the pins?

(ii) gives the stack an inner-product oracle interface: a query is a weight `W` on the cube, the
answer is `Σ_w W(w)·q(w)` (specification Definition 3.13). The probe checks that
  1. ArkLib's `OracleInterface` accepts `Weight n` (leanerVM's structure, in `Type`) as a query type;
  2. the evaluation query is the special case of the equality weight, by the built lemma
     `eqWeight_pair`;
  3. ArkLib's `Commitment.Scheme` over that interface has as opening statement
     (commitment, weight, claimed value): an inner-product commitment scheme;
  4. the opening phase "one challenge λ, then one weighted query" is an ArkLib `OracleVerifier`
     over a one-round schedule.
A wrapper `Stack` is used so that the probe's instance does not overlap the built `evalOracle`.
-/
import LeanerVM.Protocol.ClaimWeights
import LeanerVM.Protocol.ToArkLib.KeepOracles
import ArkLib.OracleReduction.Security.Basic

open LeanerVM.Protocol LeanerVM.Parameters CompPoly OracleComp OracleSpec ProtocolSpec

namespace Probe

structure Stack (n : ℕ) where
  col : Column n

/-- (ii): a query is a weight, the answer is the pairing. -/
instance innerProductOracle (n : ℕ) : OracleInterface (Stack n) where
  Query := Weight n
  toOC :=
    { spec := (Weight n) →ₒ E
      impl := fun W ↦ do return W.pair (← read).col }

theorem answer_eq (n : ℕ) (q : Stack n) (W : Weight n) :
    OracleInterface.answer q W = W.pair q.col := rfl

/-- The evaluation oracle of Layer 0 is the special case `W = eq(p, ·)`. -/
theorem answer_eqWeight (n : ℕ) (q : Stack n) (p : Vector E n) :
    OracleInterface.answer q (eqWeight p)
      = CMlPolynomialEval.eval₂Mle q.col.values (algebraMap K E) p :=
  eqWeight_pair p q.col

/-- Verbatim copy of `Commitment.Opening` (ArkLib `dca90385`,
`ArkLib/Commitments/Functional/Basic.lean:59-64`), whose module the build of `main` did not
compile (no leanerVM module imports it). -/
structure OpeningCopy {ι : Type} (oSpec : OracleSpec ι) (Data Commitment Decommitment ComKey VerifKey : Type)
    [O : OracleInterface Data] {n : ℕ} (pSpec : ProtocolSpec n) where
  opening : (ComKey × VerifKey) →
    Proof oSpec (Commitment × (q : O.Query) × O.Response q) (Data × Decommitment) pSpec

/-- Over the inner-product interface the statement of the opening is a commitment, a weight and
a claimed value: ArkLib's functional commitment scheme is then an inner-product commitment
scheme. -/
example (n : ℕ) {ι : Type} (oSpec : OracleSpec ι) (Cm Dec CK VK : Type) {m : ℕ}
    (pSpec : ProtocolSpec m)
    (S : OpeningCopy oSpec (Stack n) Cm Dec CK VK pSpec) (ck : CK) (vk : VK) :
    Proof oSpec (Cm × (W : Weight n) × E) (Stack n × Dec) pSpec :=
  S.opening (ck, vk)

/-! ## The opening phase as "λ, then one weighted query" -/

/-- One verifier message, the batching challenge. -/
@[reducible]
def openSpec : ProtocolSpec 1 := ⟨!v[.V_to_P], !v[E]⟩

instance : ∀ i, OracleInterface (openSpec.Message i)
  | ⟨0, h⟩ => nomatch h

instance : ∀ i, SampleableType (openSpec.Challenge i)
  | ⟨0, _⟩ => (inferInstance : SampleableType E)

abbrev OneStack (n : ℕ) : Fin 1 → Type := fun _ ↦ Stack n

/-- A pooled claim. -/
structure Claim (n : ℕ) where
  weight : Weight n
  value : E

/-- `W_λ = Σ_j λ^j W_j` on the cube. (The probe takes the extension as the evaluator; the closed
form the compiled verifier runs is `Σ_j λ^j W_j.mle`, equal to it by linearity.) -/
def batchWeight {n : ℕ} (lam : E) (cs : List (Claim n)) : Weight n where
  onCube := Vector.ofFn fun i ↦ (cs.zipIdx.map fun c ↦ lam ^ c.2 * c.1.weight.onCube.get i).sum
  mle := fun r ↦ CMlPolynomialEval.evalMle _ r
  mle_eq := fun _ ↦ rfl

/-- `C_λ = Σ_j λ^j c_j`. -/
def batchValue {n : ℕ} (lam : E) (cs : List (Claim n)) : E :=
  (cs.zipIdx.map fun c ↦ lam ^ c.2 * c.1.value).sum

/-- The one query of the opening phase. -/
def queryStack {n : ℕ} (W : Weight n) : OracleComp [OneStack n]ₒ E :=
  liftM <| OracleSpec.query (show [OneStack n]ₒ.Domain from ⟨0, W⟩)

/-- The opening verifier of (ii): read λ, query the stack at `W_λ`, compare with `C_λ`. -/
def openVerifier (n : ℕ) :
    OracleVerifier []ₒ (List (Claim n)) (OneStack n) Unit (OneStack n) openSpec where
  verify := fun cs chals ↦ do
    let lam : E := chals ⟨0, rfl⟩
    let a ← liftM (queryStack (batchWeight lam cs))
    if a = batchValue lam cs then pure () else failure
  outputOracle := .inl (keepOracles (OneStack n) openSpec)

#check @openVerifier
#print axioms openVerifier
#print axioms answer_eqWeight

end Probe
\end{leancode}
\src{.claude/reports/blueprint-review/probes/gt-opening-compile/InnerProductOracle.lean; OpeningCopy is a verbatim copy of ArkLib dca90385, ArkLib/Commitments/Functional/Basic.lean:59-64}

\paragraph{Command.}
\begin{shellcode}
flock .claude/reports/blueprint-review/logs/lean.lock lake env lean .claude/reports/blueprint-review/probes/gt-opening-compile/InnerProductOracle.lean
\end{shellcode}

\paragraph{Output} on 2026-09-29 (old pins: leanerVM \code{b435631}, ArkLib \code{dca90385}) and on
2026-09-30 (new pins: ArkLib \code{7653a901}, Lean 4.34.1), identical up to the date:
\begin{shellcode}
InnerProductOracle.lean:56:23: warning: Variable name `W` is not explicitly referenced. […]
openVerifier : (n : ℕ) → OracleVerifier []ₒ (List (Claim n)) (OneStack n) Unit (OneStack n) openSpec
'Probe.openVerifier' depends on axioms: [propext, Classical.choice, Quot.sound]
'Probe.answer_eqWeight' depends on axioms: [propext, Classical.choice, Quot.sound]
exit=0
\end{shellcode}
The warning is the unused binder \code{W} in the copied \lean{OpeningCopy}; harmless. The two
declarations the conclusion rests on use only the three standard axioms.

\subsection{\texttt{FiatShamirChain.lean}: what \texttt{Verifier.fiatShamir} takes, and a chain plugged in}\label{sec:gen-opening-probe-fs}

\paragraph{What it establishes.} What ArkLib's Fiat--Shamir transform takes as the challenge oracle,
and whether a concrete hash chain can be plugged in (used in \cref{sec:gen-opening-iff-compiled}).
\lean{Verifier.fiatShamir} is copied verbatim from ArkLib \code{dca90385},
\file{ArkLib/OracleReduction/FiatShamir/Basic.lean:129-136}, with the section variables of
\code{:68-70} made explicit; everything it uses (\lean{NonInteractiveVerifier},
\lean{fsChallengeOracle}, \lean{Messages.deriveTranscriptFS}) is in compiled modules.

\paragraph{Source.}
\begin{leancode}
/-
Probe (task gt-opening-compile): what does ArkLib's Fiat–Shamir transform take as the challenge
oracle, and can a concrete hash chain be plugged in?

`ArkLib/OracleReduction/FiatShamir/Basic.lean` is not compiled by the build of `main` (no leanerVM
module imports it), so `Verifier.fiatShamir` is copied verbatim from ArkLib `dca90385`,
`FiatShamir/Basic.lean:129-136`, with the section variables of `:68-70` made explicit. Everything
it uses (`NonInteractiveVerifier`, `fsChallengeOracle`, `Messages.deriveTranscriptFS`) is in
compiled modules.
-/
import ArkLib.OracleReduction.Security.Basic

open ProtocolSpec OracleComp OracleSpec

namespace Probe

variable {n : ℕ} {pSpec : ProtocolSpec n} {ι : Type} {oSpec : OracleSpec ι}
  {StmtIn StmtOut : Type}
  [VCVCompatible StmtIn] [∀ i, VCVCompatible (pSpec.Challenge i)]

/-- Verbatim: the (slow) Fiat-Shamir transformation for the verifier. -/
def fiatShamirCopy (V : Verifier oSpec StmtIn StmtOut pSpec) :
    NonInteractiveVerifier (∀ i, pSpec.Message i) (oSpec + fsChallengeOracle StmtIn pSpec)
      StmtIn StmtOut where
  verify := fun stmtIn proof => do
    let messages : pSpec.Messages := proof 0
    let transcript ← (messages.deriveTranscriptFS (oSpec := oSpec) stmtIn)
    Option.getM (← (V.verify stmtIn transcript).run)

/-- The index of a challenge-oracle query: a round that is a challenge, with the statement and
the messages sent before it. -/
example : (fsChallengeOracle StmtIn pSpec).Domain
    = ((i : pSpec.ChallengeIdx) × (StmtIn × pSpec.MessagesUpTo i.1.castSucc)) := rfl

/-- A "hash chain" in the only form the transform can see: a function from the statement and the
messages so far to the challenge of each round. -/
abbrev Chain (StmtIn : Type) (pSpec : ProtocolSpec n) :=
  (i : pSpec.ChallengeIdx) → StmtIn × pSpec.MessagesUpTo i.1.castSucc → pSpec.Challenge i

/-- The chain as a deterministic implementation of the challenge oracle (no shared oracle). -/
def chainImpl (chain : Chain StmtIn pSpec) :
    QueryImpl ([]ₒ + fsChallengeOracle StmtIn pSpec) Id
  | .inl q => PEmpty.elim q
  | .inr ⟨i, t⟩ => chain i t

/-- The compiled verifier run under a concrete chain: a deterministic function of the statement
and of ALL the prover's messages. -/
def runWithChain (V : Verifier []ₒ StmtIn StmtOut pSpec) (chain : Chain StmtIn pSpec)
    (stmt : StmtIn) (messages : ∀ i, pSpec.Message i) : Option StmtOut :=
  (simulateQ (chainImpl chain) ((fiatShamirCopy V).verify stmt (fun | 0 => messages)).run).run

#check @runWithChain
#print axioms runWithChain

-- An oracle verifier is not a `Verifier`: it must be converted first, and the conversion's
-- statement and messages are the oracles in the clear.
#check @OracleVerifier.toVerifier

end Probe
\end{leancode}
\src{.claude/reports/blueprint-review/probes/gt-opening-compile/FiatShamirChain.lean; fiatShamirCopy is a verbatim copy of ArkLib dca90385, ArkLib/OracleReduction/FiatShamir/Basic.lean:129-136}

\paragraph{Command.}
\begin{shellcode}
flock .claude/reports/blueprint-review/logs/lean.lock lake env lean .claude/reports/blueprint-review/probes/gt-opening-compile/FiatShamirChain.lean
\end{shellcode}

\paragraph{Output} (both runs, 2026-09-29 at the old pins and 2026-09-30 at the new pins):
\begin{shellcode}
@runWithChain : {n : ℕ} → {pSpec : ProtocolSpec n} → {StmtIn StmtOut : Type} →
  Verifier []ₒ StmtIn StmtOut pSpec → Chain StmtIn pSpec → StmtIn →
  ((i : pSpec.MessageIdx) → pSpec.Message i) → Option StmtOut
'Probe.runWithChain' depends on axioms: [propext, Quot.sound]
@OracleVerifier.toVerifier : … → OracleVerifier oSpec StmtIn OStmtIn StmtOut OStmtOut pSpec →
  Verifier oSpec (StmtIn × ((i : ιₛᵢ) → OStmtIn i)) (StmtOut × ((i : ιₛₒ) → OStmtOut i)) pSpec
exit=0
\end{shellcode}

\paragraph{Reading.} The copied \lean{Verifier.fiatShamir} elaborates without the
\lean{VCVCompatible} instances of its file's \code{variable} line (they are unused by the definition);
the challenge oracle's query is the pair (statement, messages so far) and any function of that shape
is an implementation; the compiled verifier under a chain is a deterministic \code{Option StmtOut}; an
oracle verifier's plain form carries the oracle statements and messages in the clear.

\subsection{\texttt{whir\_params.py} and \texttt{bus\_depth.py}: the parameters and the bus depth (scratch, Python 3; no Lean)}\label{sec:gen-opening-probe-params}

\paragraph{What \file{whir_params.py} establishes.} It ports \code{whir_config.rs:260-311, 419-707}
(\file{crates/pcs/src/whir_config.rs} at \code{a386121f}: the ladder and the per-level search) to
Python floats and compares with the pinned \file{python-verifier/verifier.py} (a copy extracted with
\code{git show a386121f:python-verifier/verifier.py} as \file{verifier_pinned.py}, since the working
tree moved). It reproduces the query table of \code{verifier.py:910} for every $(ρ, μ)$ and reads
off $η$, $m$, the Johnson list bound $L$ and the per-term security bits
(\cref{tab:gen-opening-params}).

\paragraph{Command.} In the probe directory, run on 2026-09-29 against the then-pinned working tree
and on 2026-09-30 against the extracted copy; same output:
\begin{shellcode}
PYTHONDONTWRITEBYTECODE=1 python3 whir_params.py
\end{shellcode}

\paragraph{Source.}
\begin{srccode}
# Scratch port of crates/pcs/src/whir_config.rs (pin a386121f) parameter derivation,
# to reproduce the query table of python-verifier/verifier.py:910 and to read off
# eta, m, the Johnson list bound L and the per-term security bits.
import math, sys
sys.dont_write_bytecode = True
sys.path.insert(0, ".")

SECURITY_BITS = 128
QUERY_GRINDING_BITS = 17
INITIAL_FOLDING_FACTOR = 6
SUBSEQUENT_FOLDING_FACTOR = 4
RS_INIT_RED = 3
RS_SUB_RED = 1
RESIDUAL_MAX_LOG = 5
LOG_Q = 192.0
RING_DEG = (1 << 31) + (1 << 15) + (1 << 7) + (1 << 3) + (1 << 1) + 1
M_MAX = 4096

def reduced_rate(lir, cols):
    dim = 2.0 ** cols
    return (dim - 1.0) / 2.0 ** (cols + lir)

def m_param(lir, cols, eta):
    s = math.sqrt(reduced_rate(lir, cols))
    return float(max(math.ceil(s / eta), 3))

def thm_log_a(lir, eta, cols):
    rho = reduced_rate(lir, cols)
    s = math.sqrt(rho)
    gamma = 1.0 - s - eta
    m = m_param(lir, cols, eta)
    half = m + 0.5
    num = 2.0 * half ** 5 + 3.0 * half * gamma * rho
    den = 3.0 * rho ** 1.5
    n = 2.0 ** (cols + lir)
    a = (num / den) * n + half / s
    return math.log2(a)

def log_a(lir, eta, cols, ilv):
    return thm_log_a(lir, eta, cols) + max(ilv - 1.0, 0.0)

def per_q(lir, cols, eta):
    rho = reduced_rate(lir, cols)
    gamma = 1.0 - math.sqrt(rho) - eta
    return math.log2(1.0 / (1.0 - gamma))

def list_log2(lir, cols, eta):
    rho = reduced_rate(lir, cols)
    return math.log2(1.0 / (2.0 * eta * math.sqrt(rho)))

def alg_bits(lir, cols, eta, prevq, ood):
    deg = max(RING_DEG, prevq + ood, 2)
    return LOG_Q - math.log2(deg) - list_log2(lir, cols, eta)

def ood_bits(lir, cols, eta, mu, s):
    l = list_log2(lir, cols, eta)
    lm = math.log2(mu)
    if s == 0:
        return LOG_Q - l - lm
    return s * (LOG_Q - lm) - (2.0 * l - 1.0)

def eta_for_m(lir, cols, m):
    s = math.sqrt(reduced_rate(lir, cols))
    eta = s / m
    import struct
    while m_param(lir, cols, eta) > m:
        eta = math.nextafter(eta, math.inf)
    return eta

def optimize(level, lir, cols, ilv, prevq):
    target = float(SECURITY_BITS)
    qtarget = float(max(SECURITY_BITS - QUERY_GRINDING_BITS, 1))
    mu = cols + ilv
    block = 1 << (cols + lir)
    best = None
    for m in range(3, M_MAX + 1):
        eta = eta_for_m(lir, cols, m)
        max_eta = 1.0 - math.sqrt(reduced_rate(lir, cols))
        if eta >= max_eta:
            continue
        pg = LOG_Q - log_a(lir, eta, cols, ilv)
        if pg + 1e-12 < target:
            break
        pq = per_q(lir, cols, eta)
        if not math.isfinite(pq) or pq <= 0:
            continue
        q = math.ceil(qtarget / pq)
        if q > block:
            continue
        if level == 0:
            s = 0
        else:
            s = next((s for s in range(1, 9) if ood_bits(lir, cols, eta, mu, s) + 1e-12 >= target), None)
            if s is None:
                continue
        if ood_bits(lir, cols, eta, mu, s) + 1e-12 < target or alg_bits(lir, cols, eta, prevq, s) + 1e-12 < target:
            continue
        if best is None or q < best["queries"]:
            best = dict(eta=eta, queries=q, ood=s, m=m, pg=pg, pq=pq,
                        L=2.0 ** list_log2(lir, cols, eta),
                        oodbits=ood_bits(lir, cols, eta, mu, s),
                        alg=alg_bits(lir, cols, eta, prevq, s),
                        qbits=q * pq, rho=reduced_rate(lir, cols),
                        gamma=1.0 - math.sqrt(reduced_rate(lir, cols)) - eta)
    return best

def ladder(log_n, lir):
    rates, cols, ilv, ks = [lir], [log_n - INITIAL_FOLDING_FACTOR], [INITIAL_FOLDING_FACTOR], [INITIAL_FOLDING_FACTOR]
    n_run, r_run, f_run, red = log_n - INITIAL_FOLDING_FACTOR, lir, INITIAL_FOLDING_FACTOR, RS_INIT_RED
    while n_run > RESIDUAL_MAX_LOG:
        k = min(SUBSEQUENT_FOLDING_FACTOR, n_run)
        nxt = n_run - k
        rate = r_run + (f_run - red)
        red = RS_SUB_RED
        rates.append(rate); cols.append(nxt); ilv.append(k); ks.append(k)
        n_run -= k; r_run = rate; f_run = k
    return rates, cols, ilv, ks, n_run

def derive(log_n, lir):
    rates, cols, ilv, ks, yr = ladder(log_n, lir)
    levels, prevq = [], 0
    for i in range(len(rates)):
        b = optimize(i, rates[i], cols[i], ilv[i], prevq)
        assert b is not None, (log_n, lir, i)
        b.update(rate=rates[i], cols=cols[i], ilv=ilv[i], k=ks[i])
        levels.append(b)
        prevq = b["queries"]
    return levels, yr

if __name__ == "__main__":
    import verifier_pinned as V
    ok = True
    for lir in range(1, 5):
        for mu in range(15, 29):
            levels, yr = derive(mu, lir)
            mine = tuple(l["queries"] for l in levels)
            tab = V.WHIR_QUERIES[lir - 1][mu - 15]
            cfg = V.derive_config(mu, lir)
            same = (mine == tab) and tuple(l["rate"] for l in levels) == cfg.log_inv_rates and tuple(l["k"] for l in levels) == cfg.folds
            ok &= same
            if not same:
                print("MISMATCH", lir, mu, mine, tab)
    print("query table of verifier.py:910 reproduced from the whir_config.rs derivation:", ok)
    for (mu, lir) in [(15, 1), (18, 1), (22, 1), (28, 1), (28, 2), (28, 4)]:
        levels, yr = derive(mu, lir)
        print(f"\nmu={mu} log_inv_rate={lir} residual={yr}")
        for i, l in enumerate(levels):
            print(f"  L{i}: fold k={l['k']} rate=2^-{l['rate']} msg_cols=2^{l['cols']} rho_red={l['rho']:.6f} "
                  f"m={l['m']} eta={l['eta']:.5f} gamma={l['gamma']:.5f} L<= {l['L']:.2f} "
                  f"queries={l['queries']} ood={l['ood']} | bits: mca-fold={l['pg']:.2f} query={l['qbits']:.2f}(+17 grind) "
                  f"ood={l['oodbits']:.2f} algebraic(list-unioned)={l['alg']:.2f}")
\end{srccode}
\src{.claude/reports/blueprint-review/probes/gt-opening-compile/whir_params.py}

\paragraph{Output} (2026-09-30):
\begin{shellcode}
query table of verifier.py:910 reproduced from the whir_config.rs derivation: True

mu=15 log_inv_rate=1 residual=5
  L0: fold k=6 rate=2^-1 msg_cols=2^9 rho_red=0.499023 m=395 eta=0.00179 gamma=0.29180 L<= 395.77 queries=223 ood=0 | bits: mca-fold=132.94 query=111.00(+17 grind) ood=179.46 algebraic(list-unioned)=152.37
  L1: fold k=4 rate=2^-4 msg_cols=2^5 rho_red=0.060547 m=306 eta=0.00080 gamma=0.75313 L<= 2526.97 queries=55 ood=1 | bits: mca-fold=133.22 query=111.00(+17 grind) ood=167.22 algebraic(list-unioned)=149.70

mu=18 log_inv_rate=1 residual=4
  L0: fold k=6 rate=2^-1 msg_cols=2^12 rho_red=0.499878 m=311 eta=0.00227 gamma=0.29071 L<= 311.08 queries=224 ood=0 | bits: mca-fold=131.67 query=111.00(+17 grind) ood=179.55 algebraic(list-unioned)=152.72
  L1: fold k=4 rate=2^-4 msg_cols=2^8 rho_red=0.062256 m=70 eta=0.00356 gamma=0.74692 L<= 562.20 queries=56 ood=1 | bits: mca-fold=140.88 query=111.01(+17 grind) ood=171.15 algebraic(list-unioned)=151.87
  L2: fold k=4 rate=2^-7 msg_cols=2^4 rho_red=0.007324 m=19 eta=0.00450 gamma=0.90991 L<= 1297.07 queries=32 ood=1 | bits: mca-fold=146.52 query=111.12(+17 grind) ood=169.32 algebraic(list-unioned)=150.66

mu=22 log_inv_rate=1 residual=4
  L0: fold k=6 rate=2^-1 msg_cols=2^16 rho_red=0.499992 m=216 eta=0.00327 gamma=0.28962 L<= 216.00 queries=225 ood=0 | bits: mca-fold=130.29 query=111.00(+17 grind) ood=179.79 algebraic(list-unioned)=153.25
  L1: fold k=4 rate=2^-4 msg_cols=2^12 rho_red=0.062485 m=80 eta=0.00312 gamma=0.74691 L<= 640.16 queries=56 ood=1 | bits: mca-fold=135.93 query=111.01(+17 grind) ood=170.36 algebraic(list-unioned)=151.68
  L2: fold k=4 rate=2^-7 msg_cols=2^8 rho_red=0.007782 m=42 eta=0.00210 gamma=0.90968 L<= 2698.54 queries=32 ood=1 | bits: mca-fold=137.03 query=111.00(+17 grind) ood=166.62 algebraic(list-unioned)=149.60
  L3: fold k=4 rate=2^-10 msg_cols=2^4 rho_red=0.000916 m=7 eta=0.00432 gamma=0.96542 L<= 3822.93 queries=23 ood=1 | bits: mca-fold=145.91 query=111.64(+17 grind) ood=166.20 algebraic(list-unioned)=149.10

mu=28 log_inv_rate=1 residual=2
  L0: fold k=6 rate=2^-1 msg_cols=2^22 rho_red=0.500000 m=110 eta=0.00643 gamma=0.28647 L<= 110.00 queries=228 ood=0 | bits: mca-fold=129.15 query=111.02(+17 grind) ood=180.41 algebraic(list-unioned)=154.22
  L1: fold k=4 rate=2^-4 msg_cols=2^18 rho_red=0.062500 m=81 eta=0.00309 gamma=0.74691 L<= 648.00 queries=56 ood=1 | bits: mca-fold=129.84 query=111.01(+17 grind) ood=169.86 algebraic(list-unioned)=151.66
  L2: fold k=4 rate=2^-7 msg_cols=2^14 rho_red=0.007812 m=46 eta=0.00192 gamma=0.90969 L<= 2944.18 queries=32 ood=1 | bits: mca-fold=130.39 query=111.01(+17 grind) ood=165.78 algebraic(list-unioned)=149.48
  L3: fold k=4 rate=2^-10 msg_cols=2^10 rho_red=0.000976 m=8 eta=0.00390 gamma=0.96486 L<= 4100.00 queries=23 ood=1 | bits: mca-fold=139.15 query=111.11(+17 grind) ood=165.19 algebraic(list-unioned)=149.00
  L4: fold k=4 rate=2^-13 msg_cols=2^6 rho_red=0.000120 m=4 eta=0.00274 gamma=0.98630 L<= 16644.06 queries=18 ood=1 | bits: mca-fold=140.20 query=111.41(+17 grind) ood=161.63 algebraic(list-unioned)=146.98
  L5: fold k=4 rate=2^-16 msg_cols=2^2 rho_red=0.000011 m=5 eta=0.00068 gamma=0.99594 L<= 218453.33 queries=14 ood=1 | bits: mca-fold=134.67 query=111.22(+17 grind) ood=154.94 algebraic(list-unioned)=143.26

mu=28 log_inv_rate=2 residual=2
  L0: fold k=6 rate=2^-2 msg_cols=2^22 rho_red=0.250000 m=82 eta=0.00610 gamma=0.49390 L<= 164.00 queries=113 ood=0 | bits: mca-fold=128.75 query=111.02(+17 grind) ood=179.84 algebraic(list-unioned)=153.64
  L1: fold k=4 rate=2^-5 msg_cols=2^18 rho_red=0.031250 m=43 eta=0.00411 gamma=0.81911 L<= 688.00 queries=45 ood=1 | bits: mca-fold=131.87 query=111.01(+17 grind) ood=169.69 algebraic(list-unioned)=151.57
  L2: fold k=4 rate=2^-8 msg_cols=2^14 rho_red=0.003906 m=40 eta=0.00156 gamma=0.93594 L<= 5120.31 queries=28 ood=1 | bits: mca-fold=128.89 query=111.00(+17 grind) ood=164.19 algebraic(list-unioned)=148.68
  L3: fold k=4 rate=2^-11 msg_cols=2^10 rho_red=0.000488 m=7 eta=0.00316 gamma=0.97476 L<= 7175.01 queries=21 ood=1 | bits: mca-fold=137.55 query=111.47(+17 grind) ood=163.58 algebraic(list-unioned)=148.19
  L4: fold k=4 rate=2^-14 msg_cols=2^6 rho_red=0.000060 m=3 eta=0.00258 gamma=0.98967 L<= 24966.10 queries=17 ood=1 | bits: mca-fold=139.51 query=112.14(+17 grind) ood=160.46 algebraic(list-unioned)=146.39
  L5: fold k=4 rate=2^-17 msg_cols=2^2 rho_red=0.000006 m=9 eta=0.00027 gamma=0.99734 L<= 786432.00 queries=13 ood=1 | bits: mca-fold=128.22 query=111.22(+17 grind) ood=151.25 algebraic(list-unioned)=141.42

mu=28 log_inv_rate=4 residual=2
  L0: fold k=6 rate=2^-4 msg_cols=2^22 rho_red=0.062500 m=27 eta=0.00926 gamma=0.74074 L<= 216.00 queries=57 ood=0 | bits: mca-fold=131.68 query=111.01(+17 grind) ood=179.44 algebraic(list-unioned)=153.25
  L1: fold k=4 rate=2^-7 msg_cols=2^18 rho_red=0.007812 m=11 eta=0.00804 gamma=0.90358 L<= 704.00 queries=33 ood=1 | bits: mca-fold=136.47 query=111.36(+17 grind) ood=169.62 algebraic(list-unioned)=151.54
  L2: fold k=4 rate=2^-10 msg_cols=2^14 rho_red=0.000977 m=8 eta=0.00391 gamma=0.96484 L<= 4096.25 queries=23 ood=1 | bits: mca-fold=135.15 query=111.09(+17 grind) ood=164.83 algebraic(list-unioned)=149.00
  L3: fold k=4 rate=2^-13 msg_cols=2^10 rho_red=0.000122 m=4 eta=0.00276 gamma=0.98620 L<= 16400.02 queries=18 ood=1 | bits: mca-fold=136.23 query=111.22(+17 grind) ood=161.19 algebraic(list-unioned)=147.00
  L4: fold k=4 rate=2^-16 msg_cols=2^6 rho_red=0.000015 m=3 eta=0.00129 gamma=0.99483 L<= 99864.38 queries=15 ood=1 | bits: mca-fold=134.51 query=113.94(+17 grind) ood=156.46 algebraic(list-unioned)=144.39
  L5: fold k=4 rate=2^-19 msg_cols=2^2 rho_red=0.000001 m=3 eta=0.00040 gamma=0.99841 L<= 1048576.00 queries=12 ood=1 | bits: mca-fold=130.43 query=111.51(+17 grind) ood=150.42 algebraic(list-unioned)=141.00
\end{shellcode}

A caveat on floating point: the port uses Python's \code{math} where the Rust uses \code{f64}
methods; the reproduction of every query count of the table for all 56 $(ρ, μ)$ pairs is the evidence
that the two agree where it matters. The $m$, $η$ and bit values are the port's; any last-ulp
difference would not move them visibly.

\paragraph{What \file{bus_depth.py} establishes.} It runs the pinned verifier's own
\code{build_layout} and \code{bus_layout} over a covering grid of sizes to find the maximal bus depth
under the stacking window, and at the per-log caps (\cref{sec:gen-opening-piop-error}). The dossier
records its output, not a separate command line; it imports the same extracted copy
\file{verifier_pinned.py}.

\paragraph{Source.}
\begin{srccode}
# Scratch: bus depth and blueprint error sums at admissible sizes, from the pinned Python
# verifier's own layout code (python-verifier/verifier.py at a386121f).
import sys, math, itertools
sys.dont_write_bytecode = True
sys.path.insert(0, ".")
import verifier_pinned as V

print("tables:", [(t.opcode, t.width, len(t.flushes.push), len(t.flushes.pull), len(t.count_columns), t.n_constraints) for t in V.TABLES])

def depths(log_mem, taus, kbc):
    lay = V.build_layout(range(16 * 2**kbc), log_mem, taus)
    fr = (0, lay.log_memory, lay.log_bytecode)
    push = V.bus_layout(fr, lay.push)
    pull = V.bus_layout(fr, lay.pull)
    cnt = V.bus_layout((), lay.count)
    return lay.stack_log, push.depth, pull.depth, cnt.depth

best = None
# exhaustive over a coarse but covering grid: all taus equal to t or at the floor, memory, bytecode
for log_mem in range(16, 27):
    for kbc in range(0, 29):
        for t in range(0, 27):
            for pattern in itertools.product([0, 1], repeat=6):
                taus = [t if b else 0 for b in pattern]
                taus[5] = max(taus[5], 3)
                try:
                    mu, dp, dl, dc = depths(log_mem, taus, kbc)
                except V.VerificationError:
                    continue
                if not (15 <= mu <= 28):
                    continue
                if best is None or dp > best[0]:
                    best = (dp, mu, log_mem, tuple(taus), kbc, dl, dc)
print("max bus depth found under the window mu<=28:", best)
# at the per-log caps alone (ignoring the stack window)
mu, dp, dl, dc = depths(32, [32]*6, 28)
print("per-log caps (mem 32, tau 32, kbc 28), ignoring the window: stack_log", mu, "bus depth", dp, dl, dc)
\end{srccode}
\src{.claude/reports/blueprint-review/probes/gt-opening-compile/bus_depth.py}

\paragraph{Output} (2026-09-30, same as 2026-09-29):
\begin{shellcode}
tables: [(0, 15, 5, 5, 4, 0), (1, 15, 5, 5, 4, 0), (2, 8, 3, 3, 2, 0), (3, 15, 5, 5, 4, 0), (4, 14, 5, 5, 4, 2), (5, 37, 11, 11, 10, 0)]
max bus depth found under the window mu<=28: (28, 28, 16, (0, 0, 23, 0, 23, 3), 26, 28, 26)
per-log caps (mem 32, tau 32, kbc 28), ignoring the window: stack_log 41 bus depth 38 38 37
\end{shellcode}

Per table: opcode, width, push blocks, pull blocks, count columns, constraints. The grid is not
exhaustive over all $33^6$ height vectors, so ``28'' is a lower bound on the maximum that the bus
dossier's counting argument (\file{gt-bus.md}, section D.4) confirms as the maximum.
